\chapter{Complexity — Poly Time Pairing}
%%% generated by the blueprint pipeline from TCSlib.Complexity.ClassNP.PolyTimePairing
%%%%%%%%%%%%%%%%%%%%%%%%%%%%%%%%%%%%%%%%%%%%%%%%%%%%%%%%%%%%%%%%%%%%%%%%%%%%%%%
\section{Overview}

This module proves closure properties of the class $\mathrm{FP}$ of polynomial-time
computable functions on bit strings [AB09, ch.~2] that are needed when a reduction must keep
its input while computing from it: constant functions and prepending a fixed word, the
threaded payload map on pairs, pairing of two polynomial-time functions through the
self-delimiting pair encoding $\langle a,b\rangle$, the total pair projections and their
rearrangements, concatenation, Boolean branching and conjunction, length tests on pairs, and
the unary maps $x \mapsto 1^{|x|}$ and $x \mapsto 1^{C(|x|+1)^d}$ (the input of a uniformity
machine [AB09, Def.~6.12]). All of them are assembled from a catalog of explicit finite
Turing-machine constructions.

%%%%%%%%%%%%%%%%%%%%%%%%%%%%%%%%%%%%%%%%%%%%%%%%%%%%%%%%%%%%%%%%%%%%%%%%%%%%%%%
\section{Declarations}

\begin{theorem}[Linear-time functions are polynomial-time]
\lean{Complexity.polyTimeComputable_of_linear}
\label{Complexity.polyTimeComputable_of_linear}
\leanok
\uses{Complexity.PolyTimeComputable, Turing.FinTM, Turing.FinTM.ComputesFunInTime}
\difficulty{2}
Let $f$ be a function from bit strings to bit strings. If there are a finite binary Turing machine $M$ and a constant $a \in \mathbb{N}$ such that $M$ computes $f$ within $a(n+1)$ steps on every input of length $n$, then $f$ is polynomial-time computable, i.e.\ some finite binary Turing machine computes $f$ within $C(n+1)^c$ steps on every input of length $n$, for some constants $C, c \in \mathbb{N}$.
\end{theorem}

\begin{theorem}[Constant functions are polynomial-time]
\lean{Complexity.polyTimeComputable_const}
\label{Complexity.polyTimeComputable_const}
\leanok
\uses{Complexity.PolyTimeComputable, Complexity.polyTimeComputable_of_linear, Turing.FinTM.computesFunInTime_const}
\difficulty{1}
For every fixed bit string $w$, the constant function $x \mapsto w$ on bit strings is polynomial-time computable, i.e.\ some finite binary Turing machine computes it within $C(n+1)^c$ steps on every input of length $n$, for some constants $C, c \in \mathbb{N}$.
\end{theorem}

\begin{definition}[Threaded payload map]
\lean{Complexity.pairMapSnd}
\label{Complexity.pairMapSnd}
\leanok
\uses{Turing.pairDecode, Turing.pairEncode}
Let $g$ be a function from bit strings to bit strings, and let $\langle a,b\rangle$ denote the self-delimiting pair encoding of bit strings $a$ and $b$ (every bit of $a$ doubled, then $01$, then $b$). The threaded payload map of $g$ sends a bit string $z$ to $\langle a, g(b)\rangle$ when $z = \langle a,b\rangle$ is a pair encoding, carrying the first component unchanged, and to the empty word when $z$ is not of this form.
\end{definition}

\begin{theorem}[Payload map on an encoded pair]
\lean{Complexity.pairMapSnd_pairEncode}
\label{Complexity.pairMapSnd_pairEncode}
\leanok
\uses{Complexity.pairMapSnd, Turing.pairDecode_pairEncode, Turing.pairEncode}
\difficulty{1}
Let $g$ be a function from bit strings to bit strings, let $\langle a,b\rangle$ denote the self-delimiting pair encoding of bit strings $a,b$ (every bit of $a$ doubled, then $01$, then $b$), and let $\sigma_g$ be the threaded payload map of $g$, which decodes a bit string $z$ as a pair $\langle a,b\rangle$ and outputs $\langle a, g(b)\rangle$, or outputs the empty word if $z$ is not a pair encoding. Then $\sigma_g(\langle a,b\rangle) = \langle a, g(b)\rangle$ for all bit strings $a$ and $b$.
\end{theorem}

\begin{theorem}[Payload map preserves polynomial time]
\lean{Complexity.PolyTimeComputable.pairMapSnd}
\label{Complexity.PolyTimeComputable.pairMapSnd}
\leanok
\uses{Complexity.PolyTimeComputable, Complexity.pairMapSnd, Turing.FinTM.ComputesInTime.mono, Turing.FinTM.computesFunInTime_pairMapSnd}
\difficulty{4}
Let $g$ be a polynomial-time computable function on bit strings (computed by some finite binary Turing machine within $C(n+1)^c$ steps on inputs of length $n$, for some constants $C, c$). Then its threaded payload map, which sends a pair encoding $\langle a,b\rangle$ (every bit of $a$ doubled, then $01$, then $b$) to $\langle a, g(b)\rangle$ and every bit string that is not a pair encoding to the empty word, is also polynomial-time computable.
\end{theorem}

\begin{lemma}[Appending to a pair encoding]
\lean{Complexity.pairEncode_append}
\label{Complexity.pairEncode_append}
\leanok
\uses{Turing.pairEncode}
\difficulty{1}
Let $\langle a,b\rangle$ denote the self-delimiting pair encoding of bit strings $a$ and $b$ (every bit of $a$ doubled, then $01$, then $b$). For all bit strings $a, b, c$, appending $c$ to $\langle a,b\rangle$ appends it to the second component: $\langle a,b\rangle\, c = \langle a, b\,c\rangle$, where juxtaposition denotes concatenation.
\end{lemma}

\begin{theorem}[Pairing of polynomial-time functions]
\lean{Complexity.PolyTimeComputable.pairEncode}
\label{Complexity.PolyTimeComputable.pairEncode}
\leanok
\uses{Complexity.PolyTimeComputable, Complexity.PolyTimeComputable.comp, Complexity.PolyTimeComputable.pairMapSnd, Complexity.pairEncode_append, Complexity.pairMapSnd_pairEncode, Complexity.polyTimeComputable_const, Complexity.polyTimeComputable_of_linear, Turing.FinTM.computesFunInTime_pairConcat, Turing.FinTM.computesFunInTime_pairDup, Turing.FinTM.computesFunInTime_pairFst, Turing.FinTM.computesFunInTime_pairSnd, Turing.pairDecode_pairEncode, Turing.pairEncode}
\difficulty{5}
Let $f$ and $g$ be polynomial-time computable functions on bit strings (each computed by some finite binary Turing machine within $C(n+1)^c$ steps on inputs of length $n$, for some constants $C, c$). Then the map $x \mapsto \langle f(x), g(x)\rangle$ is polynomial-time computable, where $\langle a,b\rangle$ denotes the self-delimiting pair encoding of bit strings $a$ and $b$ (every bit of $a$ doubled, then $01$, then $b$).
\end{theorem}

\begin{lemma}[Stripping the last one of an all-ones word]
\lean{Complexity.splitAtLastTrue_replicate_succ}
\label{Complexity.splitAtLastTrue_replicate_succ}
\leanok
\uses{Turing.splitAtLastTrue}
\difficulty{3}
For a bit string $v$, let $\mathrm{split}(v)$ be the prefix of $v$ preceding its last occurrence of $1$, undefined when $v$ contains no $1$. Then for every $n \in \mathbb{N}$, $\mathrm{split}(1^{n+1})$ is defined and equals $1^n$, where $1^k$ is the word of $k$ ones.
\end{lemma}

\begin{theorem}[Unary length map is polynomial-time]
\lean{Complexity.polyTimeComputable_unary}
\label{Complexity.polyTimeComputable_unary}
\leanok
\uses{Complexity.PolyTimeComputable, Complexity.PolyTimeComputable.comp, Complexity.PolyTimeComputable.pairEncode, Complexity.polyTimeComputable_id, Complexity.polyTimeComputable_of_linear, Complexity.splitAtLastTrue_replicate_succ, Turing.FinTM.computesFunInTime_pairSnd, Turing.FinTM.computesFunInTime_polyUnary, Turing.FinTM.computesFunInTime_stripLast, Turing.pairDecode, Turing.pairDecode_pairEncode, Turing.pairEncode, Turing.splitAtLastTrue}
\difficulty{4}
The map sending a bit string $x$ to the word $1^{|x|}$ of $|x|$ ones is polynomial-time computable, i.e.\ some finite binary Turing machine computes it within $C(n+1)^c$ steps on every input of length $n$, for some constants $C, c \in \mathbb{N}$.
\end{theorem}

\begin{theorem}[Closure of FP under concatenation]
\lean{Complexity.PolyTimeComputable.append}
\label{Complexity.PolyTimeComputable.append}
\leanok
\uses{Complexity.PolyTimeComputable, Complexity.PolyTimeComputable.comp, Complexity.PolyTimeComputable.pairEncode, Complexity.polyTimeComputable_of_linear, Turing.FinTM.computesFunInTime_pairConcat, Turing.pairDecode_pairEncode}
\difficulty{3}
Let $f$ and $g$ be polynomial-time computable functions on bit strings (each computed by some finite binary Turing machine within $C(n+1)^c$ steps on inputs of length $n$, for some constants $C, c$). Then the map $x \mapsto f(x)\,g(x)$, sending $x$ to the concatenation of $f(x)$ and $g(x)$, is polynomial-time computable.
\end{theorem}

\begin{theorem}[Polynomial unary generator is polynomial-time]
\lean{Complexity.polyTimeComputable_polyUnary}
\label{Complexity.polyTimeComputable_polyUnary}
\leanok
\uses{Complexity.PolyTimeComputable, Turing.FinTM.computesFunInTime_polyUnary}
\difficulty{2}
For all $C, d \in \mathbb{N}$, the map sending a bit string $x$ to the word $1^{C(|x|+1)^d}$ of $C(|x|+1)^d$ ones is polynomial-time computable, i.e.\ some finite binary Turing machine computes it within $K(n+1)^k$ steps on every input of length $n$, for some constants $K, k \in \mathbb{N}$.
\end{theorem}

\begin{theorem}[Prepending a fixed word is polynomial-time]
\lean{Complexity.polyTimeComputable_prepend}
\label{Complexity.polyTimeComputable_prepend}
\leanok
\uses{Complexity.PolyTimeComputable, Complexity.polyTimeComputable_of_linear, Turing.FinTM.computesFunInTime_prepend}
\difficulty{1}
For every fixed bit string $w$, the map $x \mapsto w\,x$ that prepends $w$ to a bit string $x$ is polynomial-time computable, i.e.\ some finite binary Turing machine computes it within $C(n+1)^c$ steps on every input of length $n$, for some constants $C, c \in \mathbb{N}$.
\end{theorem}

\begin{definition}[Total first projection of a pair]
\lean{Complexity.pairFstD}
\label{Complexity.pairFstD}
\leanok
\uses{Turing.pairDecode}
Let $\langle a,b\rangle$ denote the self-delimiting pair encoding of bit strings $a$ and $b$ (every bit of $a$ doubled, then $01$, then $b$). The total first projection $\pi_1$ sends a bit string $z$ to $a$ when $z = \langle a,b\rangle$ is a pair encoding, and to the empty word when $z$ is not of this form.
\end{definition}

\begin{definition}[Total second projection of a pair]
\lean{Complexity.pairSndD}
\label{Complexity.pairSndD}
\leanok
\uses{Turing.pairDecode}
Let $\langle a,b\rangle$ denote the self-delimiting pair encoding of bit strings $a$ and $b$ (every bit of $a$ doubled, then $01$, then $b$). The total second projection $\pi_2$ sends a bit string $z$ to $b$ when $z = \langle a,b\rangle$ is a pair encoding, and to the empty word when $z$ is not of this form.
\end{definition}

\begin{theorem}[First projection of an encoded pair]
\lean{Complexity.pairFstD_pairEncode}
\label{Complexity.pairFstD_pairEncode}
\leanok
\uses{Complexity.pairFstD, Turing.pairDecode_pairEncode, Turing.pairEncode}
\difficulty{1}
Let $\langle a,b\rangle$ denote the self-delimiting pair encoding of bit strings $a$ and $b$ (every bit of $a$ doubled, then $01$, then $b$), and let $\pi_1$ be the total first projection, which returns the first component of a pair encoding and the empty word on any other bit string. Then $\pi_1(\langle a,b\rangle) = a$ for all bit strings $a, b$.
\end{theorem}

\begin{theorem}[Second projection of an encoded pair]
\lean{Complexity.pairSndD_pairEncode}
\label{Complexity.pairSndD_pairEncode}
\leanok
\uses{Complexity.pairSndD, Turing.pairDecode_pairEncode, Turing.pairEncode}
\difficulty{1}
Let $\langle a,b\rangle$ denote the self-delimiting pair encoding of bit strings $a$ and $b$ (every bit of $a$ doubled, then $01$, then $b$), and let $\pi_2$ be the total second projection, which returns the second component of a pair encoding and the empty word on any other bit string. Then $\pi_2(\langle a,b\rangle) = b$ for all bit strings $a, b$.
\end{theorem}

\begin{theorem}[First projection does not increase length]
\lean{Complexity.length_pairFstD_le}
\label{Complexity.length_pairFstD_le}
\leanok
\uses{Complexity.pairFstD, Turing.eq_pairEncode_of_pairDecode, Turing.length_pairEncode, Turing.pairDecode}
\difficulty{4}
Let $\pi_1$ be the total first projection, which sends a bit string of the form $\langle a,b\rangle$ (the self-delimiting pair encoding: every bit of $a$ doubled, then $01$, then $b$) to $a$, and any other bit string to the empty word. Then for every bit string $z$, $|\pi_1(z)| \le |z|$.
\end{theorem}

\begin{theorem}[First projection is polynomial-time]
\lean{Complexity.polyTimeComputable_pairFstD}
\label{Complexity.polyTimeComputable_pairFstD}
\leanok
\uses{Complexity.PolyTimeComputable, Complexity.pairFstD, Complexity.polyTimeComputable_of_linear, Turing.FinTM.computesFunInTime_pairFst}
\difficulty{1}
The total first projection $\pi_1$, which sends a bit string of the form $\langle a,b\rangle$ (the self-delimiting pair encoding: every bit of $a$ doubled, then $01$, then $b$) to $a$ and any other bit string to the empty word, is polynomial-time computable, i.e.\ some finite binary Turing machine computes it within $C(n+1)^c$ steps on every input of length $n$, for some constants $C, c \in \mathbb{N}$.
\end{theorem}

\begin{theorem}[Second projection is polynomial-time]
\lean{Complexity.polyTimeComputable_pairSndD}
\label{Complexity.polyTimeComputable_pairSndD}
\leanok
\uses{Complexity.PolyTimeComputable, Complexity.pairSndD, Complexity.polyTimeComputable_of_linear, Turing.FinTM.computesFunInTime_pairSnd}
\difficulty{1}
The total second projection $\pi_2$, which sends a bit string of the form $\langle a,b\rangle$ (the self-delimiting pair encoding: every bit of $a$ doubled, then $01$, then $b$) to $b$ and any other bit string to the empty word, is polynomial-time computable, i.e.\ some finite binary Turing machine computes it within $C(n+1)^c$ steps on every input of length $n$, for some constants $C, c \in \mathbb{N}$.
\end{theorem}

\begin{theorem}[Iterated first projection is polynomial-time]
\lean{Complexity.polyTimeComputable_iterate_pairFstD}
\label{Complexity.polyTimeComputable_iterate_pairFstD}
\leanok
\uses{Complexity.PolyTimeComputable, Complexity.PolyTimeComputable.comp, Complexity.pairFstD, Complexity.polyTimeComputable_id, Complexity.polyTimeComputable_pairFstD}
\difficulty{4}
Let $\pi_1$ be the total first projection, which sends a bit string of the form $\langle a,b\rangle$ (the self-delimiting pair encoding: every bit of $a$ doubled, then $01$, then $b$) to $a$ and any other bit string to the empty word. For every $n \in \mathbb{N}$, the $n$-fold iterate $\pi_1^{n} = \pi_1 \circ \cdots \circ \pi_1$ (the identity for $n = 0$), which extracts the root of a left-nested tuple, is polynomial-time computable.
\end{theorem}

\begin{theorem}[Concatenating pair components is polynomial-time]
\lean{Complexity.polyTimeComputable_pairConcat}
\label{Complexity.polyTimeComputable_pairConcat}
\leanok
\uses{Complexity.PolyTimeComputable, Complexity.pairFstD, Complexity.pairSndD, Complexity.polyTimeComputable_of_linear, Turing.FinTM.computesFunInTime_pairConcat, Turing.pairDecode}
\difficulty{4}
Let $\pi_1$ and $\pi_2$ be the total projections, which send a bit string of the form $\langle a,b\rangle$ (the self-delimiting pair encoding: every bit of $a$ doubled, then $01$, then $b$) to $a$ and to $b$ respectively, and any other bit string to the empty word. Then the map $z \mapsto \pi_1(z)\,\pi_2(z)$, sending $z$ to the concatenation of its two components, is polynomial-time computable.
\end{theorem}

\begin{theorem}[Swapping pair components is polynomial-time]
\lean{Complexity.polyTimeComputable_pairSwap}
\label{Complexity.polyTimeComputable_pairSwap}
\leanok
\uses{Complexity.PolyTimeComputable, Complexity.PolyTimeComputable.pairEncode, Complexity.pairFstD, Complexity.pairSndD, Complexity.polyTimeComputable_pairFstD, Complexity.polyTimeComputable_pairSndD, Turing.pairEncode}
\difficulty{1}
Let $\langle a,b\rangle$ denote the self-delimiting pair encoding of bit strings $a$ and $b$ (every bit of $a$ doubled, then $01$, then $b$), and let $\pi_1$ and $\pi_2$ be the total projections, which send $\langle a,b\rangle$ to $a$ and to $b$ respectively, and any bit string that is not a pair encoding to the empty word. Then the map $z \mapsto \langle \pi_2(z), \pi_1(z)\rangle$ is polynomial-time computable.
\end{theorem}

\begin{theorem}[Polynomial-time branching]
\lean{Complexity.polyTimeComputable_ite}
\label{Complexity.polyTimeComputable_ite}
\leanok
\uses{Complexity.PolyTimeComputable, Turing.FinTM.ComputesInTime.mono, Turing.FinTM.computesFunInTime_cond}
\difficulty{5}
Let $p$ be a Boolean-valued predicate on bit strings and let $f, g$ be functions from bit strings to bit strings. Suppose that the map sending $x$ to the one-bit word $p(x)$ is polynomial-time computable, and that $f$ and $g$ are polynomial-time computable (each computed by some finite binary Turing machine within $C(n+1)^c$ steps on inputs of length $n$, for some constants $C, c$). Then the map sending $x$ to $f(x)$ if $p(x)$ is true and to $g(x)$ otherwise is polynomial-time computable.
\end{theorem}

\begin{theorem}[Polynomial-time conjunction of one-bit tests]
\lean{Complexity.polyTimeComputable_and}
\label{Complexity.polyTimeComputable_and}
\leanok
\uses{Complexity.PolyTimeComputable, Complexity.polyTimeComputable_const, Complexity.polyTimeComputable_ite}
\difficulty{3}
Let $p$ and $q$ be Boolean-valued predicates on bit strings such that the maps sending $x$ to the one-bit words $p(x)$ and $q(x)$ are both polynomial-time computable (computed by some finite binary Turing machine within $C(n+1)^c$ steps on inputs of length $n$, for some constants $C, c$). Then the map sending $x$ to the one-bit word $p(x) \land q(x)$ is polynomial-time computable.
\end{theorem}

\begin{theorem}[Polynomial-time length comparison on pairs]
\lean{Complexity.polyTimeComputable_lenLe}
\label{Complexity.polyTimeComputable_lenLe}
\leanok
\uses{Complexity.PolyTimeComputable, Complexity.PolyTimeComputable.comp, Complexity.PolyTimeComputable.pairEncode, Complexity.pairFstD, Complexity.pairSndD, Complexity.polyTimeComputable_pairFstD, Complexity.polyTimeComputable_pairSndD, Complexity.polyTimeComputable_prepend, Turing.FinTM.computesFunInTime_pairLenCheck, Turing.pairDecode, Turing.pairDecode_pairEncode, Turing.pairEncode}
\difficulty{4}
Let $\pi_1$ and $\pi_2$ be the total projections, which send a bit string of the form $\langle a,b\rangle$ (the self-delimiting pair encoding: every bit of $a$ doubled, then $01$, then $b$) to $a$ and to $b$ respectively, and any other bit string to the empty word. Then the map sending a bit string $z$ to the one-bit word that is $1$ if $|\pi_2(z)| \le |\pi_1(z)|$ and $0$ otherwise is polynomial-time computable, i.e.\ some finite binary Turing machine computes it within $C(n+1)^c$ steps on every input of length $n$, for some constants $C, c \in \mathbb{N}$.
\end{theorem}

\begin{theorem}[Polynomial-time length equality on pairs]
\lean{Complexity.polyTimeComputable_lenEq}
\label{Complexity.polyTimeComputable_lenEq}
\leanok
\uses{Complexity.PolyTimeComputable, Complexity.PolyTimeComputable.comp, Complexity.pairFstD, Complexity.pairFstD_pairEncode, Complexity.pairSndD, Complexity.pairSndD_pairEncode, Complexity.polyTimeComputable_and, Complexity.polyTimeComputable_lenLe, Complexity.polyTimeComputable_pairSwap}
\difficulty{4}
Let $\pi_1$ and $\pi_2$ be the total projections, which send a bit string of the form $\langle a,b\rangle$ (the self-delimiting pair encoding: every bit of $a$ doubled, then $01$, then $b$) to $a$ and to $b$ respectively, and any other bit string to the empty word. Then the map sending a bit string $z$ to the one-bit word that is $1$ if $|\pi_1(z)| = |\pi_2(z)|$ and $0$ otherwise is polynomial-time computable, i.e.\ some finite binary Turing machine computes it within $C(n+1)^c$ steps on every input of length $n$, for some constants $C, c \in \mathbb{N}$.
\end{theorem}
